\chapter{OFDM}\label{ch:ofdm}

Orthogonal frequency-division multiplexing turns a frequency-selective channel into many
parallel flat channels, each equalized by a single complex division. This chapter derives the
property, describes the subcarrier layout and the pilot-based receiver of
\file{phy/waveform/ofdm.py}, and examines the two weaknesses of OFDM: synchronization and
peak-to-average power ratio.

% ------------------------------------------------------------------
\section{Cyclic prefix and circular convolution}
An OFDM symbol carries $N$ frequency-domain values $X_k$ (the subcarrier symbols). The
transmitter computes the unitary inverse DFT and prepends the last $N_{\mathrm{cp}}$ samples
(the \emph{cyclic prefix}):
\begin{equation}
  x[n]=\frac{1}{\sqrt N}\sum_{k=0}^{N-1}X_k\,e^{j2\pi kn/N},\qquad n=-N_{\mathrm{cp}},\dots,N-1 .
\end{equation}
\begin{proposition}
If the channel impulse response has length $L\le N_{\mathrm{cp}}+1$, then after removing the
prefix and applying the unitary DFT,
\begin{equation}
  Y_k=H_kX_k+W_k,\qquad H_k=\sum_{l=0}^{L-1}h_l\,e^{-j2\pi kl/N},\qquad W_k\sim\CN(0,\sigma^2)\ \text{i.i.d.}
  \label{eq:ofdm-model}
\end{equation}
\end{proposition}
\begin{proof}
Thanks to the prefix, the $N$ retained received samples equal the \emph{circular} convolution of
$x$ with $h$: $\vect y=\mat C\vect x+\vect w$, with $\mat C$ circulant. Every circulant matrix is
diagonalized by the DFT, $\mat C=\mat F\Her\diag(H_0,\dots,H_{N-1})\mat F$ with $\mat F$ the
unitary DFT matrix. Hence $\mat F\vect y=\diag(H)\mat F\vect x+\mat F\vect w$, and
$\mat F\vect w$ is white with the same variance because $\mat F$ is unitary.
\end{proof}

\begin{figure}[htbp]
\centering
\begin{tikzpicture}[x=1cm,y=1cm]
  \foreach \s in {0,1,2} {
    \begin{scope}[xshift=\s*4.3cm]
      \draw[fill=kgrid!35,draw=kgrid] (0,0) rectangle (0.9,0.8);
      \draw[fill=kgrid!10,draw=kgrid] (0.9,0) rectangle (4.1,0.8);
      \fill[kgrid!35] (3.2,0) rectangle (4.1,0.8);
      \node[lbl] at (0.45,0.4) {CP};
      \node[lbl] at (2.05,0.4) {$N$ samples};
      \draw[->,kgrid!80!black,thin] (3.65,0.85) to[bend right=30] (0.45,0.85);
    \end{scope}}
  \draw[decorate,decoration={brace,mirror,amplitude=4pt}] (0,-0.1) -- node[below=4pt,lbl]{$N+N_{\mathrm{cp}}$} (4.1,-0.1);
  \draw[accent2,thick] (5.2,-0.55) -- (8.4,-0.55) node[right,lbl]{FFT window};
  \draw[accent2] (5.2,-0.45)--(5.2,-0.65) (8.4,-0.45)--(8.4,-0.65);
\end{tikzpicture}
\caption{OFDM symbols with cyclic prefix. Any FFT window starting inside the prefix (after the
channel transient) sees a cyclically shifted symbol, which only multiplies $X_k$ by a linear phase.}
\label{fig:ofdm-cp}
\end{figure}

\begin{incode}
\blk{ofdm_modulator} builds an \code{OFDMLayout}, maps data and pilots onto the grid, applies
\code{np.fft.ifft(..., norm="ortho")} and adds the prefix. The unitary transform keeps the
noise variance per subcarrier equal to the noise variance per time sample.

It sets \meta{sps}$=N/N_{\mathrm{act}}$ and
\meta{power}$=(N_{\mathrm{data}}P_d+N_{\mathrm{pil}})/N$. By \cref{eq:esn0} the AWGN block in
$E_s/N_0$ mode then produces $\sigma^2=1/(E_s/N_0)$, the per-subcarrier SNR. Cyclic-prefix and
pilot overheads are deliberately not counted in $E_s$.
\end{incode}
\codeexcerpt{ofdm_mod}{Resource mapping, unitary IFFT and cyclic prefix.}

% ------------------------------------------------------------------
\section{Subcarrier layout}
$N_{\mathrm{act}}$ subcarriers are used around DC. DC itself is left empty, as in Wi-Fi and LTE,
and the remaining edge subcarriers act as guard bands. Every \prm{pilot_spacing}-th active
subcarrier carries a known pilot. The last active subcarrier is always a pilot, so that
interpolation never extrapolates. Pilot values are a fixed pseudo-random QPSK sequence;
identical pilots would add up coherently into a periodic time-domain peak.

\begin{figure}[htbp]
\centering
\begin{tikzpicture}[x=0.36cm,y=1cm]
  % active k = -13..-1, 1..13; pilots where the active index is a multiple of 6, plus the last one
  \foreach \k in {-16,-15,-14,14,15,16} \draw[black!25] (\k,0)--(\k,0.25);
  \draw[black!60] (0,0)--(0,0.25); \node[lbl,below] at (0,0) {DC};
  \foreach \k in {-12,-11,-10,-9,-8,-6,-5,-4,-3,-2,1,2,3,4,5,7,8,9,10,11}
    \fill[accent3!70] (\k-0.3,0) rectangle (\k+0.3,0.65);
  \foreach \k in {-13,-7,-1,6,12,13}
    \fill[accent2] (\k-0.3,0) rectangle (\k+0.3,0.95);
  \draw[->] (-17,0) -- (17.5,0) node[right,lbl]{$k$};
  \fill[accent2] (-15,1.45) rectangle (-14.4,1.7); \node[lbl,right] at (-14.2,1.57) {pilot};
  \fill[accent3!70] (-9,1.45) rectangle (-8.4,1.7); \node[lbl,right] at (-8.2,1.57) {data};
  \draw[black!40] (-3,1.45)--(-3,1.7); \node[lbl,right] at (-2.8,1.57) {unused (guard)};
\end{tikzpicture}
\caption{Layout with $N_{\mathrm{act}}=26$ active subcarriers and a pilot every 6 (the first and the
last active subcarriers are pilots).}
\label{fig:ofdm-layout}
\end{figure}

% ------------------------------------------------------------------
\section{Channel estimation and equalization}
\subsection{Least squares on pilots}
At pilot subcarriers \cref{eq:ofdm-model} gives the LS estimate $\hat H_p=Y_p/X_p$, with error
variance $\sigma^2$ since pilots have unit modulus. The channel at data subcarriers is
obtained by linear interpolation, implemented as a fixed matrix from pilot to active
subcarriers. Optionally the estimates are averaged over the OFDM symbols of a frame, which
lowers the estimation noise when the channel is static.

\subsection{Removing the timing phase slope}
A timing offset of $\delta$ samples inside the prefix multiplies $H_k$ by $e^{-j2\pi k\delta/N}$.
This fast phase rotation across subcarriers is poorly approximated by linear interpolation of
the real and imaginary parts. The estimator therefore:
\begin{enumerate}
  \item estimates the dominant slope from products of adjacent pilots,
  $\hat s=\arg\sum_p\hat H_{p+1}\hat H_p\conj\,/\,\Delta_p$;
  \item removes it;
  \item interpolates the now slowly varying response;
  \item restores the slope.
\end{enumerate}
\codeexcerpt{ls_slope}{LS estimation with phase-slope removal (\blk{ofdm_channel_estimator}).}

\begin{figure}[htbp]
\centering
\begin{tikzpicture}
\begin{groupplot}[group style={group size=1 by 2,vertical sep=6mm,x descriptions at=edge bottom},
  width=12cm,height=4.4cm,xmin=-51,xmax=51,xlabel={subcarrier index $k$}]
\nextgroupplot[ylabel={$|H_k|$ [dB]},legend pos=south west,legend columns=3]
  \addplot[th=black!70] table[x=k,y=true_db]{\datafile{ofdm_chest}};
  \addplot[th=accent,dashed] table[x=k,y=ls_db]{\datafile{ofdm_chest}};
  \addplot[sim=accent2,mark size=1.4pt] table[x=k,y=db]{\datafile{ofdm_pilots}};
  \legend{true channel,LS + interpolation,pilot estimates}
\nextgroupplot[ylabel={$\arg H_k$ [rad]}]
  \addplot[th=black!70] table[x=k,y=true_ph]{\datafile{ofdm_chest}};
  \addplot[th=accent,dashed] table[x=k,y=ls_ph]{\datafile{ofdm_chest}};
\end{groupplot}
\end{tikzpicture}
\caption{One OFDM symbol over the EVA channel ($N=128$, 100 active subcarriers, pilot every 6,
$15$\,kHz spacing, $E_s/N_0=15$\,dB): true response from the channel's \code{csi} (\prm{perfect}
estimator) and LS estimate.}
\label{fig:ofdm-chest}
\end{figure}

\subsection{One-tap equalization}
\blk{ofdm_equalizer} extracts the data subcarriers and divides by the estimate,
$\hat X_k=Y_k/\hat H_k$. It writes the per-subcarrier noise variance $\sigma^2/|\hat H_k|^2$ to
\meta{noise_var}. With coding and interleaving across subcarriers, the soft demapper then gives
low weight to faded subcarriers. This is how OFDM obtains frequency diversity (see the Wi-Fi and
LTE scenarios).

% ------------------------------------------------------------------
\section{Blind timing and frequency synchronization}\label{sec:ofdm-sync}
The cyclic prefix is a copy of the end of each symbol, so the received signal is correlated
with itself $N$ samples later over the prefix. Van de Beek \emph{et al.}~\cite{vandebeek1997}
use this redundancy. With a CFO of $\varepsilon$ subcarrier spacings, $r[k+N]\approx r[k]e^{j2\pi\varepsilon}$
inside the prefix, so the correlation
\begin{equation}
  \gamma(\theta)=\sum_{m}\sum_{k=\theta+mL}^{\theta+mL+N_{\mathrm{cp}}-1}r[k]\,r\conj[k+N],\qquad L=N+N_{\mathrm{cp}},
\end{equation}
peaks in magnitude at the symbol start, and its phase gives the fractional CFO:
\begin{equation}
  \hat\theta=\argmax_\theta|\gamma(\theta)|,\qquad \hat\varepsilon=-\frac{1}{2\pi}\arg\gamma(\hat\theta),\quad|\varepsilon|<\tfrac12 .
\end{equation}
The block accumulates $\gamma$ over all OFDM symbols of the frame with cumulative sums, which
makes the search over $\theta$ linear in the frame length. It then starts the FFT window
\prm{backoff} samples early, inside the prefix, to stay clear of the channel's delay spread.
The resulting linear phase is absorbed by the channel estimator. Compared with the full ML
metric of van de Beek, the energy correction term is omitted; it matters only at low SNR.
\codeexcerpt{ofdm_sync}{CP correlation, timing and fractional CFO (\blk{ofdm_sync}).}

\begin{pitfall}
An integer CFO (a multiple of the subcarrier spacing) shifts the whole grid and is
invisible to the CP correlation. It must be resolved with a preamble or pilots, which this
block does not do.
\end{pitfall}

% ------------------------------------------------------------------
\section{Peak-to-average power ratio}\label{sec:papr}
Each time sample is a sum of $N_{\mathrm{act}}$ independent terms. By the central limit theorem
it is nearly $\CN(0,P)$, so $|x|^2/P$ is approximately exponential. Treating the $N$ Nyquist-rate
samples of a symbol as independent gives the classic approximation
\begin{equation}
  P\big(\mathrm{PAPR}>\gamma\big)\approx1-\big(1-e^{-\gamma}\big)^{N_{\mathrm{act}}} .
\end{equation}

\begin{figure}[htbp]
\centering
\begin{tikzpicture}
\begin{semilogyaxis}[width=10cm,height=6cm,xlabel={$\gamma$ [dB]},ylabel={$P(\mathrm{PAPR}>\gamma)$},
  xmin=4,xmax=13,ymin=1e-3,ymax=1,legend pos=south west]
  \addplot[sim=accent] table[x=papr,y=sim64]{\datafile{papr_ccdf}};
  \addplot[th=accent] table[x=papr,y=th64]{\datafile{papr_ccdf}};
  \addplot[sim={accent3,mark=square*}] table[x=papr,y=sim256]{\datafile{papr_ccdf}};
  \addplot[th=accent3] table[x=papr,y=th256]{\datafile{papr_ccdf}};
  \addplot[sim={accent2,mark=triangle*}] table[x=papr,y=sim1024]{\datafile{papr_ccdf}};
  \addplot[th=accent2] table[x=papr,y=th1024]{\datafile{papr_ccdf}};
  \legend{$N{=}64$ sim.,approx.,$N{=}256$ sim.,approx.,$N{=}1024$ sim.,approx.}
\end{semilogyaxis}
\end{tikzpicture}
\caption{CCDF of the per-symbol PAPR of QPSK-OFDM at the output of \blk{ofdm_modulator} (no
oversampling, no prefix) and the approximation above. A PAPR of 10--12\,dB forces the
power-amplifier back-off discussed in \cref{ch:impairments}.}
\label{fig:papr}
\end{figure}

\begin{remark}
The modulator produces Nyquist-rate samples, which miss part of the analog peaks; with
oversampling the PAPR is typically about 0.5\,dB higher.
\end{remark}
